% =====================================================================
%  1bat-ud1-12.tex  —  HORA 12: Tocant pressupostos gegants (operacions i interpretació)
% =====================================================================
\horatitol{Sessió 12 — Tocant pressupostos gegants: operacions i interpretació}%
{Operar amb notació científica sobre xifres reals i valorar si l'ordre de magnitud del resultat té sentit.}

\subsection{Idea clau}

Ja sabem escriure xifres grans en notació científica. Avui les fem servir per
\textbf{operar}: multiplicar i dividir pressupostos i poblacions, i sobretot
\textbf{interpretar} si el resultat té sentit a la vida real.

\begin{idea}
\textbf{Operar en notació científica:}
\begin{itemize}[leftmargin=1.4em, itemsep=2pt]
  \item \textbf{Multiplicar}: es multipliquen les mantisses i se \textbf{sumen} els
        exponents. \quad $(a\times 10^{m})\cdot(b\times 10^{n})=(a\cdot b)\times 10^{m+n}$.
  \item \textbf{Dividir}: es divideixen les mantisses i es \textbf{resten} els
        exponents. \quad $\dfrac{a\times 10^{m}}{b\times 10^{n}}=\dfrac{a}{b}\times 10^{m-n}$.
  \item Si la mantissa del resultat se'n surt de $[1,10)$, es \textbf{reajusta}.
\end{itemize}
\textbf{Pregunta clau sempre:} «aquest resultat \emph{té sentit} a la vida real?».
L'ordre de magnitud (l'exponent) ens ho diu de seguida.
\end{idea}

\subsection{Exemples}

\begin{exemple}[euros per persona]
Repartim el pressupost $3,326\times 10^{9}\eur$ entre $1,95\times 10^{6}$ persones en
risc. Dividim mantisses i restem exponents:
\[
\frac{3,326\times 10^{9}}{1,95\times 10^{6}}
=\frac{3,326}{1,95}\times 10^{9-6}
\approx 1,71\times 10^{3}\eur.
\]
És a dir, uns $\num{1710}\eur$ per persona. \textbf{Té sentit?} Sí: és una xifra de
l'ordre dels milers d'euros, raonable per a un ajut anual.
\end{exemple}

\begin{exemple}[una multiplicació]
Si una ajuda mensual de $4\times 10^{2}\eur$ (és a dir $400\eur$) es paga a
$5\times 10^{4}$ persones, el cost total mensual és:
\[
(4\times 10^{2})\cdot(5\times 10^{4})=(4\cdot 5)\times 10^{2+4}
=20\times 10^{6}=2\times 10^{7}\eur.
\]
(Hem reajustat $20\times 10^{6}$ a $2\times 10^{7}$.) Són $\num{20000000}\eur$, és a
dir $20$ milions: un ordre de magnitud creïble per a un programa d'aquesta mida.
\end{exemple}

\subsection{Exercicis resolts}

\begin{exresolt}[dividir dues xifres grans]
Una partida de $1,8\times 10^{8}\eur$ es reparteix entre $3\times 10^{5}$
beneficiaris. Quants euros toquen per persona?
\solucio
Dividim mantisses i restem exponents:
\[
\frac{1,8\times 10^{8}}{3\times 10^{5}}=\frac{1,8}{3}\times 10^{8-5}
=0,6\times 10^{3}.
\]
Reajustem ($0,6\times 10^{3}=6\times 10^{2}$): $600\eur$ per persona.
\resposta{$6\times 10^{2}=600\eur$ per persona.}
\end{exresolt}

\begin{exresolt}[valorar l'ordre de magnitud]
Algú calcula que, repartint un pressupost de $3\times 10^{9}\eur$ entre $2\times 10^{6}$
persones, toquen $1,5\times 10^{6}\eur$ per persona. És plausible? Troba l'error.
\solucio
Fem la divisió correctament: $\dfrac{3\times 10^{9}}{2\times 10^{6}}=1,5\times 10^{3}$,
és a dir $\num{1500}\eur$ per persona. El resultat d'aquella persona, $1,5\times 10^{6}$
($\num{1500000}\eur$ per persona!), és \textbf{absurd}: ningú no rep un milió i mig.
L'error està en l'\textbf{exponent}: va restar malament ($9-6=3$, no $6$).
\resposta{No és plausible; el correcte és $1,5\times 10^{3}=\num{1500}\eur$ per persona.}
\end{exresolt}

\subsection{Exercicis per fer}

\begin{exercicis}
\begin{exlist}
  \item Multiplica (i reajusta si cal):
        \begin{enumerate}[label=\alph*), leftmargin=1.6em, topsep=2pt]
          \item $(2\times 10^{3})\cdot(3\times 10^{4})$
          \item $(5\times 10^{2})\cdot(6\times 10^{5})$
        \end{enumerate}
  \item Divideix (i reajusta si cal):
        \begin{enumerate}[label=\alph*), leftmargin=1.6em, topsep=2pt]
          \item $\dfrac{8\times 10^{7}}{4\times 10^{3}}$
          \item $\dfrac{2,4\times 10^{6}}{6\times 10^{2}}$
        \end{enumerate}
  \item Un pressupost de $3,326\times 10^{9}\eur$ es reparteix entre $1,95\times 10^{6}$
        persones. Calcula els euros per persona (com a l'exemple) i digues si l'ordre
        de magnitud té sentit.
  \item Una ajuda de $300\eur$ mensuals es paga a $2\times 10^{4}$ famílies. Escriu
        $300$ en notació científica, calcula el cost mensual total i expressa'l en
        notació científica.
  \item \textbf{(Ampliació)} Compara dues partides: una de $4\times 10^{8}\eur$ per a
        $5\times 10^{5}$ persones i una altra de $9\times 10^{8}\eur$ per a
        $1,5\times 10^{6}$ persones. Quina dóna més euros per persona?
  \item \textbf{(Sentit del resultat)} En repartir un pressupost municipal entre els
        habitants, surten $45\eur$ per habitant. Si en repartir el pressupost de tot un
        país (molt més gran) entre molta més gent surten $1\,200\eur$ per habitant,
        és coherent que la xifra per habitant sigui més gran? Comenta-ho.
\end{exlist}
\end{exercicis}

% ---------------------------------------------------------------------
\fullsolucions{Sessió 12}
\begin{solucions}
\begin{sollist}
  \item (a) $6\times 10^{7}$; \quad (b) $30\times 10^{7}=3\times 10^{8}$.
  \item (a) $2\times 10^{4}$; \quad (b) $0,4\times 10^{4}=4\times 10^{3}$.
  \item $\dfrac{3,326}{1,95}\times 10^{3}\approx 1,71\times 10^{3}\eur\approx\num{1710}\eur$
        per persona. Sí, té sentit (ordre dels milers d'euros).
  \item $300=3\times 10^{2}$; cost $=(3\times 10^{2})\cdot(2\times 10^{4})=6\times 10^{6}\eur$
        ($6$ milions).
  \item Primera: $\dfrac{4\times 10^{8}}{5\times 10^{5}}=0,8\times 10^{3}=800\eur$/persona.
        Segona: $\dfrac{9\times 10^{8}}{1,5\times 10^{6}}=6\times 10^{2}=600\eur$/persona.
        La \textbf{primera} dóna més per persona ($800>600$).
  \item No és automàtic que un pressupost més gran doni més per habitant: depèn de la
        \emph{relació} entre pressupost i població. Aquí, $1\,200\eur$ per habitant pot
        ser coherent si el pressupost creix més de pressa que la població; cal mirar la
        divisió, no només la mida del pressupost.
\end{sollist}
\end{solucions}
